% Source: https://artofproblemsolving.com/community/c3h1271474_wait_its_different_right
% Stable identity: aops-topic-1271474-q01
The parallelogram has base $b$ and height $s\sin\theta$, so its area is $bs\sin\theta$.

\begin{asy}
usepackage("amsmath");
size(5.5cm);
defaultpen(fontsize(10pt));
draw(box((0,0),(3,2)),blue); label("$s$",(0,1),W,blue); label("$b$",(1.5,0),S,blue);
pair P=(4,0),Q=(7,0),R=Q+2*dir(55),Sq=P+2*dir(55); draw(P--Q--R--Sq--cycle,red); draw(Sq--(Sq.x,0),gray+dashed); label("$s\sin\theta$",(Sq.x,Sq.y/2),E,gray); label("$\theta$",P+(0.35,0.12),fontsize(8));
label("\mbox{}",(min(currentpicture,user=true)+max(currentpicture,user=true))/2,black); // keeps the text after the figure black when LaTeX typesets the labels inline
\end{asy}

Since $\sin\theta<1$ for $\theta<90^\circ$, the area $bs\sin\theta$ is \emph{smaller} than the rectangle's area $bs$. So the area does not stay the same; keeping the side lengths does not keep the height.
